%=============================================================================!
% 2.7  DRAG AND TERMINAL VELOCITY                                             !
%=============================================================================!
\section{Drag and terminal velocity}
\label{sec:ne-drag}

Drop a small bead into a jar of honey or a thick oil. It speeds up for a
moment and then sinks at a steady speed, however tall the jar. This is
the first force we meet that depends on the velocity, and solving its
equation of motion needs a new function, the exponential. The same
equation returns in Chapter~13, where it describes how a thermostat
steadies the temperature of a simulation.

\subsection*{The drag force}

The liquid resists the motion of the bead with a drag force that points
against the velocity. For small bodies moving slowly through a thick
liquid this force is, to a good approximation, proportional to the
velocity,
\begin{equation*}
   \vb{F}_{\mathrm{drag}} = -b\,\vb{v} ,
\end{equation*}
where the drag coefficient $b$, in kilograms per second, depends on the
size of the body and the thickness of the liquid. Fast bodies moving
through air feel a drag that grows more nearly as the square of the
speed, which we do not treat. The liquid also pushes the bead upwards
with a constant force called buoyancy, equal to the weight of the liquid
that the bead displaces. For a bead much denser than the liquid this is
a small correction, and we neglect it.

\subsection*{The equation of motion}

We measure the depth $x$ downwards, so that the weight $mg$ is positive
and the drag $-bv$ opposes a downward velocity. The second law for the
bead is
\begin{equation*}
   m\,\frac{\dd v}{\dd t} = mg - bv ,
\end{equation*}
and dividing both sides by $m$,
\begin{equation}
   \frac{\dd v}{\dd t} = g - \gamma v, \qquad \gamma = \frac{b}{m} .
   \label{eq:ne-drag}
\end{equation}
The constant $\gamma$, called the drag rate, has units of
$(\si{\kilogram\per\second})/\si{\kilogram} = \si{\per\second}$, an
inverse time; we shall see that $1/\gamma$ sets the time scale on which the
bead settles into its steady motion.

Before solving \cref{eq:ne-drag} we can find the steady speed directly.
The bead stops speeding up when $\dd v/\dd t = 0$, which happens when the
drag has grown to balance the weight,
\begin{equation*}
   0 = g - \gamma v_\infty
   \quad\Longrightarrow\quad
   v_\infty = \frac{g}{\gamma} = \frac{mg}{b} .
\end{equation*}
This steady speed is called the terminal velocity. We write it
$v_\infty$, read `v infinity', because the bead approaches it as the time
grows without limit.

\subsection*{Measuring the velocity from its terminal value}

\Cref{eq:ne-drag} becomes simpler if we follow not the velocity itself
but its distance from the terminal velocity, $w = v - v_\infty$. Since
$v_\infty$ is a constant, $\dd w/\dd t = \dd v/\dd t$, and substituting
$v = w + v_\infty$ into \cref{eq:ne-drag},
\begin{align*}
   \frac{\dd w}{\dd t}
   &= g - \gamma\,(w + v_\infty)
      && \text{substituting for $v$}\\
   &= (g - \gamma v_\infty) - \gamma w
      && \text{multiplying out the bracket}\\
   &= -\gamma w
      && \text{since $\gamma v_\infty = g$.}
\end{align*}
The gap between the velocity and its terminal value therefore shrinks at
a rate proportional to the gap itself: when the gap is large it closes
quickly, and as it narrows it closes ever more slowly. To solve this
equation we need a function whose derivative is a constant times the
function, and the boxes below build it.

\begin{toolbox}[The exponential function]
Define $\mathrm{E}(s)$ as the function that equals its own derivative and
takes the value 1 at $s = 0$,
\begin{equation*}
   \frac{\dd\mathrm{E}}{\dd s} = \mathrm{E}(s), \qquad \mathrm{E}(0) = 1 .
\end{equation*}
Differentiating both sides again shows that the second derivative also
equals $\mathrm{E}$, and so on for every derivative, so at $s = 0$ every
derivative of $\mathrm{E}$ equals 1. The Taylor series of
\cref{sec:mo-taylor} about $s = 0$ is therefore
\begin{equation*}
   \mathrm{E}(s) = 1 + s + \frac{s^2}{2!} + \frac{s^3}{3!} + \dotsb .
\end{equation*}
Differentiating the series term by term confirms that it is its own
derivative: since $n! = n\times(n-1)!$, the term $s^n/n!$ becomes
$n\,s^{n-1}/n! = s^{n-1}/(n-1)!$, the term before it, so the
differentiated series is the same series again. Such a function
therefore exists, and since $\dd\mathrm{E}/\dd s = \mathrm{E}$ is a
first-order differential equation, the uniqueness theorem of
\cref{sec:ne-equations} allows only one with $\mathrm{E}(0) = 1$.
Its value at $s = 1$ is a number called $\mathrm{e}$. Adding the first
seven terms, $1 + 1 + 0.5 + 0.1667 + 0.0417 + 0.0083 + 0.0014 = 2.7181$;
the next three terms, $0.0002$, $0.00002$ and $0.000003$, raise this to
$2.71828$, and later terms are too small to change these digits:
$\mathrm{e} = 2.71828\ldots$.

Now let $\lambda$ be any constant. By the chain rule of
\cref{sec:mo-acceleration}, with inner function $\lambda t$,
\begin{equation*}
   \frac{\dd}{\dd t}\,\mathrm{E}(\lambda t)
   = \mathrm{E}'(\lambda t)\,\lambda = \lambda\,\mathrm{E}(\lambda t) ,
\end{equation*}
the second step using $\mathrm{E}' = \mathrm{E}$,
so $\mathrm{E}(\lambda t)$ is a function whose derivative is $\lambda$
times itself, which is what we need. The series also gives values at
negative $s$, for instance
\begin{equation*}
   \mathrm{E}(-1) = 1 - 1 + 0.5 - 0.1667 + 0.0417 - 0.0083 + 0.0014 - \dotsb
   = 0.368 .
\end{equation*}
\end{toolbox}

\begin{toolbox}[The product rule]
For two functions $p(s)$ and $q(s)$, the change in their product over a
small step $h$ can be split into a change in $p$ and a change in $q$, by
adding and subtracting $p(s)\,q(s+h)$,
\begin{equation*}
   p(s+h)\,q(s+h) - p(s)\,q(s)
   = \bigl[p(s+h) - p(s)\bigr]\,q(s+h) + p(s)\,\bigl[q(s+h) - q(s)\bigr] .
\end{equation*}
Dividing by $h$ and letting $h \to 0$, the brackets divided by $h$ become
the derivatives $p'$ and $q'$, and $q(s+h)$ becomes $q(s)$, so
\begin{equation*}
   \frac{\dd}{\dd s}\bigl(pq\bigr) = p'q + pq' .
\end{equation*}
\end{toolbox}

The product rule shows why $\mathrm{E}$ deserves its usual name. Fix a
number $c$ and consider $f(s) = \mathrm{E}(c + s)\,\mathrm{E}(-s)$. The
derivative of $\mathrm{E}(-s)$ is $-\mathrm{E}(-s)$, by the chain rule
with $\lambda = -1$, and the derivative of $\mathrm{E}(c + s)$ is
$\mathrm{E}(c + s)$, by the chain rule with inner function $c + s$, whose
derivative is 1, so
\begin{align*}
   f'(s) &= \mathrm{E}(c+s)\,\mathrm{E}(-s)
      + \mathrm{E}(c+s)\,\bigl[-\mathrm{E}(-s)\bigr]
      && \text{by the product rule}\\
   &= 0 .
\end{align*}
A function whose derivative is zero everywhere does not change, by
\cref{eq:mo-integrals} (with $f$ in place of $v$ and $f' = 0$ in place
of $a$, it reads $f(s) = f(0) + \int_0^s 0\,\dd s' = f(0)$), so $f$ keeps its value at $s = 0$, which is
$\mathrm{E}(c)\,\mathrm{E}(0) = \mathrm{E}(c)$. Setting $c = 0$ gives
$\mathrm{E}(s)\,\mathrm{E}(-s) = 1$, and dividing $\mathrm{E}(c+s)\,
\mathrm{E}(-s) = \mathrm{E}(c)$ by $\mathrm{E}(-s)$, which is the same as
multiplying by $\mathrm{E}(s)$,
\begin{equation*}
   \mathrm{E}(c + s) = \mathrm{E}(c)\,\mathrm{E}(s) .
\end{equation*}
This is the rule obeyed by powers, $\mathrm{e}^{c+s} =
\mathrm{e}^c\,\mathrm{e}^s$. Applying it with $c = s = 1$ gives
$\mathrm{E}(2) = \mathrm{e}^2$, then $\mathrm{E}(3) = \mathrm{e}^3$, and
so on, and $\mathrm{E}(-s) = 1/\mathrm{E}(s)$ matches $\mathrm{e}^{-s} =
1/\mathrm{e}^s$. Two further facts follow. Since $\mathrm{E}(s)\,
\mathrm{E}(-s) = 1$, $\mathrm{E}$ is never zero, and since $\mathrm{E}(s)
= \mathrm{E}(\tfrac s2)^2$, by the rule with both $c$ and $s$ set to
$\tfrac s2$, it is never negative; so $\mathrm{E}(s) > 0$ for every $s$.
Its derivative, equal to itself, is then positive too, so $\mathrm{E}$
rises steadily and never takes the same value at two different $s$.
Fractions agree with powers as well: $\mathrm{E}(\tfrac12)^2 =
\mathrm{E}(\tfrac12 + \tfrac12) = \mathrm{E}(1) = \mathrm{e}$, so
$\mathrm{E}(\tfrac12) = \sqrt{\mathrm{e}} = \mathrm{e}^{1/2}$, the positive
root. The
function $\mathrm{E}$ therefore extends the powers of $\mathrm{e}$ to
every real number, and from now on we write it $\mathrm{e}^s$, the
exponential function. Because it never takes a value twice, it can be
undone: its inverse is the natural logarithm, $\ln$, defined by
$\ln(\mathrm{e}^s) = s$. It is not the base-10 logarithm $\log$ of
\cref{sec:mo-taylor}; calculators mark it `ln'.

The rule for sums also explains how a decaying exponential decays.
Lowering $s$ by one multiplies the function by a fixed factor,
$\mathrm{e}^{s-1} = \mathrm{e}^s\,\mathrm{e}^{-1} = 0.368\,\mathrm{e}^s$.
In $\mathrm{e}^{\lambda t}$ with negative $\lambda$ the exponent $s =
\lambda t$ falls by $\lambda\times(1/\lvert\lambda\rvert) = -1$ each time
$t$ grows by $1/\lvert\lambda\rvert$, so the function falls to about
\SI{37}{\percent} of its value over every such interval, wherever it
starts.

\subsection*{Solving the equation}

The gap obeys $\dd w/\dd t = -\gamma w$, and by the toolbox, with
$\lambda = -\gamma$, the function $w = C\,\mathrm{e}^{-\gamma t}$ has
derivative $-\gamma C\,\mathrm{e}^{-\gamma t} = -\gamma w$ for any
constant $C$. At $t = 0$ the exponential is 1, so $C$ is the starting gap,
$C = v_0 - v_\infty$, where $v_0$ is the velocity at release. Here the
force depends on the velocity alone, so the second of
\cref{eq:ne-first-order} involves $v$ alone, and by the uniqueness
theorem of \cref{sec:ne-equations} the starting velocity fixes its
solution. We have therefore found it, and adding $v_\infty$ back,
\begin{equation}
   v(t) = v_\infty + (v_0 - v_\infty)\,\mathrm{e}^{-\gamma t} .
   \label{eq:ne-drag-v}
\end{equation}
Whatever the starting velocity, the gap shrinks to $\mathrm{e}^{-1} =
\SI{37}{\percent}$ of its size every interval $1/\gamma$, and the
velocity approaches $v_\infty$ from below if the bead starts slower and
from above if it starts faster. A bead released from rest has $v_0 = 0$,
and then $v(t) = v_\infty\,(1 - \mathrm{e}^{-\gamma t})$.

As an example, take a bead of mass \SI{0.01}{\kilogram} with a drag
coefficient of \SI{0.05}{\kilogram\per\second}. Then $\gamma = 0.05/0.01
= \SI{5}{\per\second}$, so $1/\gamma = \SI{0.2}{\second}$, and the
terminal velocity is $v_\infty = 9.80665/5 = \SI{1.961}{\metre\per
\second}$. These values are chosen so that the approach is easy to see
on a scale of seconds. A real bead this heavy and this fast would feel a
drag that is not proportional to its speed, while for the small, slow
beads to which the linear law applies the whole approach is over in a
small fraction of a second. Released from rest, the bead's velocity as a
fraction of $v_\infty$ is $1 - \mathrm{e}^{-\gamma t}$, where
$\mathrm{e}^{-2} = 0.368^2$, $\mathrm{e}^{-3} = 0.368^3$ and
$\mathrm{e}^{-5} = 0.368^5$ by the rule for sums:
\begin{center}
\begin{tabular}{cccc}
\toprule
$t$ / \si{\second} & $\gamma t$ & $\mathrm{e}^{-\gamma t}$
   & $v/v_\infty = 1 - \mathrm{e}^{-\gamma t}$\\
\midrule
0.2 & 1 & 0.368 & \SI{63.2}{\percent}\\
0.4 & 2 & 0.135 & \SI{86.5}{\percent}\\
0.6 & 3 & 0.050 & \SI{95.0}{\percent}\\
1.0 & 5 & 0.007 & \SI{99.3}{\percent}\\
\bottomrule
\end{tabular}
\end{center}
To find when it reaches \SI{90}{\percent} of $v_\infty$, we need
$\mathrm{e}^{-\gamma t} = 0.1$. Taking the natural logarithm of both
sides gives $-\gamma t = \ln 0.1 = -2.303$, the value of the logarithm
coming from a calculator, so $t = 2.303/5 =
\SI{0.461}{\second}$ (\cref{fig:ne-drag}a).

\subsection*{The distance fallen}

The depth follows by integrating the velocity, as in
\cref{sec:mo-integrating}. The function $-\mathrm{e}^{-\gamma t'}/\gamma$
has derivative $\mathrm{e}^{-\gamma t'}$, by the toolbox with $\lambda =
-\gamma$, so
\begin{equation*}
   \int_0^t \mathrm{e}^{-\gamma t'}\,\dd t'
   = \Bigl[-\frac{\mathrm{e}^{-\gamma t'}}{\gamma}\Bigr]_0^t
   = \frac{1 - \mathrm{e}^{-\gamma t}}{\gamma} ,
\end{equation*}
and for a bead released from rest at $x = 0$,
\begin{equation}
   x(t) = \int_0^t v_\infty\bigl(1 - \mathrm{e}^{-\gamma t'}\bigr)\,\dd t'
   = v_\infty\,t - \frac{v_\infty}{\gamma}\,\bigl(1 -
   \mathrm{e}^{-\gamma t}\bigr) .
   \label{eq:ne-drag-x}
\end{equation}
Two limits check this result. Long after release the exponential has
died away and $x \approx v_\infty t - v_\infty/\gamma$: the bead sinks at
its terminal velocity, but a fixed distance $v_\infty/\gamma$ behind
where it would be had it moved at that velocity from the start, because
it spent its first moments speeding up. For the example bead this lag is
$1.961/5 = \SI{0.392}{\metre}$. Just after release, when $\gamma t$ is
small, the exponential series with $s = -\gamma t$ gives
$\mathrm{e}^{-\gamma t} = 1 - \gamma t + \tfrac12\gamma^2 t^2 +
\order{t^3}$, so $1 - \mathrm{e}^{-\gamma t} = \gamma t - \tfrac12\gamma^2
t^2 + \order{t^3}$, with $\order{t^3}$ as in
\cref{sec:mo-taylor}, and substituting into
\cref{eq:ne-drag-x},
\begin{align*}
   x(t) &= v_\infty t - \frac{v_\infty}{\gamma}\Bigl(\gamma t
      - \tfrac12\gamma^2 t^2\Bigr) + \order{t^3}\\
   &= v_\infty t - v_\infty t + \tfrac12 v_\infty\gamma\, t^2 + \order{t^3}
      && \text{multiplying out the bracket}\\
   &= \tfrac12 g t^2 + \order{t^3}
      && \text{since $v_\infty\gamma = g$.}
\end{align*}
At first the bead is moving too slowly for the drag to matter, and it
falls freely. For the example bead, after \SI{0.02}{\second} it has
fallen \SI{1.898}{\milli\metre} against \SI{1.961}{\milli\metre} in free
fall; after \SI{0.1}{\second}, \SI{41.79}{\milli\metre} against
\SI{49.03}{\milli\metre}; and after a second, \SI{1.572}{\metre} against
\SI{4.903}{\metre} (\cref{fig:ne-drag}b).

\begin{figure}[htb]
\centering
\includegraphics{ch02_newton/drag}
\caption{Falling from rest with a drag proportional to the velocity. (a)
Velocity for three drag rates $\gamma$, approaching the terminal velocity
$v_\infty = g/\gamma$ (dotted), against free fall (dashed); dots mark $t
= 1/\gamma$, where \SI{63}{\percent} of the terminal velocity is reached.
(b) Distance fallen for $\gamma = \SI{5}{\per\second}$, which follows
free fall at first and then the line $v_\infty(t - 1/\gamma)$.}
\label{fig:ne-drag}
\end{figure}

\notebook{02}{change the mass and drag coefficient and watch the approach
to terminal velocity}

\begin{selfcheck}
A bead is fired downwards into the honey faster than its terminal
velocity. Does it speed up or slow down? Sketch its velocity against
time.
\end{selfcheck}
